% Fichier complété par tools/complete_missing_corrections.py.
% Source : TEC/TEC-04-Meq-Jeq/36_VisEcrou/36_VisEcrou.tex
% Statut : rédigé (4 question(s)).
\begin{corrigebox}[Corrigé rédigé]
\CorrigeQuestion{1}
Chaîne : moteur 1 $(C_m,\omega_1)$ $\rightarrow$ transmission par
                courroie/poulies $(\omega_2)$ $\rightarrow$ vis--écrou de pas $p_v$
                $\rightarrow$ piston 3 $(F_3,v_3)$.
\CorrigeQuestion{2}
Sans glissement, $\omega_2/\omega_1=D_1/D_2=1/2$. Pour la vis--écrou,
                $v_3=p_v\omega_2/(2\pi)$. Ainsi
                \[\boxed{v_3=\frac{p_v}{4\pi}\,\omega_1}.\]
\CorrigeQuestion{3}
Avec $v_3=p_v/(4\pi)\omega_1$, l'énergie $\tfrac12m_3v_3^2$ équivaut à
                $\tfrac12J_{eq}\omega_1^2$. Donc
                $J_{eq}=J_1+J_2(\omega_2/\omega_1)^2+m_3(p_v/(4\pi))^2$.
\CorrigeQuestion{4}
Ramenée à la translation :
                $M_{eq}=m_3+J_1(\omega_1/v_3)^2+J_2(\omega_2/v_3)^2$.
\end{corrigebox}
